Let \(K\succeq0\) have rank \(r\ge1\), let \(Q\succeq0\), and set \(A=K^{1/2}QK^{1/2}\). Then \(A\succeq0\), \(\operatorname{range}(A)\subseteq\operatorname{range}(K)\), and
\[
\lambda_{\max}(A)\ge \frac{\operatorname{tr}(A)}r. \tag{1}
\]
Indeed, if \(\xi_1,\ldots,\xi_r\ge0\) are the eigenvalues of \(A\) on \(\operatorname{range}(K)\), including zeros, then
\[
\lambda_{\max}(A)=\max_{1\le h\le r}\xi_h
\ge \frac1r\sum_{h=1}^r\xi_h
=\frac{\operatorname{tr}(A)}r.
\]
Equality holds if and only if \(\xi_1=\cdots=\xi_r=\operatorname{tr}(A)/r\), equivalently
\[
A=\frac{\operatorname{tr}(A)}rP_{\operatorname{range}(K)}. \tag{2}
\]
Consequently, at any proposed common design, (2) for every active class is necessary and sufficient for simultaneous equality in the classwise normalized trace inequalities at that design. For the solved rank-\(q\) block menu, the condition is \(C_j=\{\operatorname{tr}(C_j)/q\}I_q\), which the regime-specific independent-sign product laws satisfy. What remains open is the joint feasibility and optimal loss for a general noncommuting menu.